% 原书分章 PDF 第 2—7 页；章首导言与第 14.1 节。
\chapter{随机处理策略}
\label{ch:stochastic-treatment-regimes}
\chapterauthors{Iván Díaz 与 Mark J. van der Laan}
\chapteraffiliations{%
I. Díaz（通讯作者）：美国康奈尔大学威尔康奈尔医学院卫生保健政策与研究系生物统计学和流行病学部，402 East 67th Street, New York, NY 10065, USA；电子邮箱：\href{mailto:ild2005@med.cornell.edu}{ild2005@med.cornell.edu}。\\
M. J. van der Laan：美国加利福尼亚大学伯克利分校生物统计学部及统计学系，101 Haviland Hall, \#7358, Berkeley, CA 94720, USA；电子邮箱：\href{mailto:laan@berkeley.edu}{laan@berkeley.edu}。\\
本章 DOI：\url{https://doi.org/10.1007/978-3-319-65304-4_14}。}

研究因果关系的标准统计方法，将一组特定处理下的反事实结局定义为：在一个对所有个体都实施某种给定处理策略的假想世界中，所观测到的结局。例如，若处理有两个可能取值，我们可以比较分别以概率 1 分配各处理水平的策略下反事实结局的期望，以此定义因果效应。这类策略通常称为静态处理策略（static treatment regime）。另一类值得研究的策略，将个体的处理水平设定为其已测量历史的函数。这类策略称为动态处理策略（dynamic treatment regime），因为它们可以随个体已观测到的处理前特征而变化。静态策略和动态策略常被统称为确定性处理策略（deterministic treatment regime），因为它们完全由处理前测量的变量决定。

尽管确定性策略在应用研究中十分常见，但对于无法直接干预的现象，它们并不能为有关因果关系的问题提供适当框架。例如，在公共卫生研究中，现实可行的策略往往无法将处理变量置于某个确定状态（如按照一个确定性函数来设定个体的锻炼方案并不现实）；另一些策略则源于政策的实施，而政策旨在使人群行为发生随机变化。例如，通过大众媒体发布倡导使用安全套的信息，在社区层面是确定性的，在个体层面却是随机的，因为每个个体会依据外生因素决定是否接受处理。此外，确定性策略的因果效应可能无法识别，因为所关注的策略得不到观测数据的支持。例如，我们可以预期，健康问题会使一部分人无法进行较高水平的身体活动。这也使基于确定性策略的因果效应难以解释，因为估计出的效应对应的是实际中无法实施的策略。

\begin{sourcebox}
本章考虑动态策略的一种推广：允许所分配的处理还依赖处理的自然取值（natural value of treatment），即未施加干预时观测到的处理取值。由于只有在观测到处理的自然取值后，才能确定这种策略所分配的处理水平，我们将其称为随机处理策略（stochastic treatment regime）。文献中还使用随机政策、随机干预、随机化动态策略、修正处理策略（modified treatment policy）等名称。随机策略允许依赖处理的自然取值，因此总可以将它们定义为具有实际意义且现实可行的策略。
\end{sourcebox}

为说明这一点，考虑下面两个例子。

\begin{bookresult}{例}{14.1}
Tager 等（1998）开展了一项研究，主要目的是评估休闲时间身体活动（leisure-time physical activity，LTPA）对老年人死亡率的影响。原则上，我们可以考虑一组对应于 LTPA 确定性策略的假想世界，例如，将 LTPA 水平确定性地设为它的每一个可能取值。尽管理论上可以这样设想，但如此定义的反事实结局并不令人满意，因为我们不可能实施一项将个体身体活动水平确定性地设为某个值的策略。一种解决方法是考虑这样的策略：将所分配的处理设为处理自然取值的函数。例如，对当前身体活动水平为 $a$ 的个体，在该策略下可以分配 $a+\delta$。更现实的做法是，仅将这种策略分配给能够实行它的个体；是否可行可根据健康状况、当前身体活动水平等其他协变量确定。这类策略比任何确定性策略都更现实，也可以认为，它能够在现实世界中实施。Díaz 和 van der Laan（2012）最早研究了如此定义的因果效应的定义、识别与估计，Haneuse 和 Rotnitzky（2013）随后作了进一步研究。
\end{bookresult}

\begin{bookresult}{例}{14.2}
Mann 等（2010）开展了一项研究，分析空气污染水平对哮喘儿童呼吸道症状的因果效应，即弗雷斯诺哮喘儿童环境研究（Fresno Asthmatic Children's Environment Study，FACES）。该研究的一个核心目标，是考察空气中 $\mathrm{NO}_2$ 浓度对哮喘儿童喘鸣的影响。具体而言，我们可能希望评估这样一种策略的效应：降低 $\mathrm{NO}_2$ 浓度分布右尾的空气污染浓度。这类策略的一个例子，是强制将 $\mathrm{NO}_2$ 水平控制在某个阈值以下。在这一策略下，已经达标的单位，即污染水平低于所选阈值的单位，可能没有动力继续降低污染水平，因此在策略实施后可能保持不变。超过阈值的单位可能只将污染水平降低到足以达到阈值、符合政策要求的程度，却没有动力进一步降低。Díaz 和 van der Laan（2013a）发展了评估这类策略效应的因果推断方法。
\end{bookresult}

随机处理策略除了有助于定义更有实际意义的目标因果参数，还为连续暴露提供了更易处理的估计框架。因果剂量—反应曲线（causal dose-response curve）等参数在非参数模型中难以估计，因为它们不具路径可微性，例如，其效率界为无穷大。另一些方法，如将连续处理分组的方法，没有利用处理的连续性，因而也不足以回答有关干预的问题。随机处理策略还可以用来处理一些常见问题，例如自然直接效应（natural direct effect，NDE）的识别与估计、社区干预、个体化处理策略，以及意向性治疗规则。

\section{数据、记号与感兴趣的参数}

设 $A$ 表示处理变量，$Y$ 表示连续或二元结局，$W$ 表示观测到的处理前协变量向量。令 $O=(W,A,Y)$ 为服从分布 $P_0$ 的随机变量，$O_1,\ldots,O_n$ 为 $O$ 的 $n$ 个独立同分布观测所构成的样本。我们假定 $P_0\in\mathcal M$，其中 $\mathcal M$ 是非参数模型，包含 $O$ 上相对于控制测度 $\nu$ 的所有连续密度。以 $p_0$ 表示相应的概率密度函数，则
\[
 p_0(o)=p_0(y\mid a,w)p_0(a\mid w)p_0(w).
\]
记 $g_0(a\mid w)=p_0(a\mid w)$、$\bar Q_0(a,w)=E_0(Y\mid A=a,W=w)$，以及 $q_{W,0}(w)=p_0(w)$。对给定函数 $f(o)$ 和一般分布函数 $P\in\mathcal M$，还记 $Pf=\int f(o)\,dP(o)$。我们用 $P_n$ 表示 $O_1,\ldots,O_n$ 的经验分布。

假定下列非参数结构方程模型（nonparametric structural equation model，NPSEM）成立：
\begin{equation}
 W=f_W(U_W);\qquad A=f_A(W,U_A);\qquad Y=f_Y(A,W,U_Y).
 \tag{14.1}
\end{equation}
这组方程表示一个假定生成观测数据 $O$ 的机制模型，其中包含若干假设。首先，模型隐含了时间顺序：假定 $Y$ 发生在 $A$ 和 $W$ 之后，$A$ 发生在 $W$ 之后。其次，假定每个变量由位于它之前的观测变量，再加上一个记为 $U$ 的外生变量，通过确定性函数生成。每个外生变量都被假定包含相应观测变量的全部未观测原因。我们对外生变量作出以下独立性假设：
\begin{equation}
 U_A\indep U_Y.
 \tag{14.2}
\end{equation}
这一假设对于从观测数据分布识别 $A$ 对 $Y$ 的因果效应至关重要，第 14.1.1 节将说明相应的识别结果。图 14.1 给出了这一假设所允许的有向无环图（directed acyclic graph，DAG）。

我们如下定义 $A$ 对 $Y$ 的因果效应。考虑对 NPSEM (14.1) 作一项假想修改：删除与 $A$ 对应的方程，将 $A$ 设为某个假想策略 $d(A,W)$。策略 $d$ 既依赖在没有该策略时本来会分配的处理水平 $A$，也依赖协变量 $W$。在上述说明性例子中，这些策略可如下定义。

\begin{figure}[htbp]
 \centering
 % 图 14.1：每面板 6 条实线有向边，另有一条无箭头虚线相关弧。
\begin{tikzpicture}[causal,x=1.25cm,y=1.25cm,every node/.style={inner sep=2pt,font=\normalsize}]
 \begin{scope}
  \node (aUW) at (0,2.65) {$U_W$};
  \node (aW) at (0,1.7) {$W$};
  \node (aA) at (-.75,1) {$A$};
  \node (aY) at (.75,1) {$Y$};
  \node (aUA) at (-1.5,.25) {$U_A$};
  \node (aUY) at (1.5,.25) {$U_Y$};
  \draw[ed] (aUW)--(aW);
  \draw[ed] (aW)--(aA);
  \draw[ed] (aW)--(aY);
  \draw[ed] (aA)--(aY);
  \draw[ed] (aUA)--(aA);
  \draw[ed] (aUY)--(aY);
  \draw[dashed] (aUW.east) .. controls (1.42,2.02) and (1.5,1.25) .. (aUY.north);
 \end{scope}
 \begin{scope}[xshift=6.9cm]
  \node (bUW) at (0,2.65) {$U_W$};
  \node (bW) at (0,1.7) {$W$};
  \node (bA) at (-.75,1) {$A$};
  \node (bY) at (.75,1) {$Y$};
  \node (bUA) at (-1.5,.25) {$U_A$};
  \node (bUY) at (1.5,.25) {$U_Y$};
  \draw[ed] (bUW)--(bW);
  \draw[ed] (bW)--(bA);
  \draw[ed] (bW)--(bY);
  \draw[ed] (bA)--(bY);
  \draw[ed] (bUA)--(bA);
  \draw[ed] (bUY)--(bY);
  \draw[dashed] (bUW.west) .. controls (-1.42,2.02) and (-1.5,1.25) .. (bUA.north);
 \end{scope}
\end{tikzpicture}

 \caption{所允许的有向无环图。虚线表示相关关系，实线箭头表示因果关系。}
 \label{fig:ch14-allowed-dags}
\end{figure}

\begin{bookresult}{例}{14.1（续）}
设给定 $W=w$ 时 $A$ 的条件分布支撑于区间 $(l(w),u(w))$。也就是说，协变量为 $W=w$ 的个体所能达到的最大身体活动量为 $u(w)$。于是可以定义
\begin{equation}
 d(a,w)=
 \begin{cases}
 a+\delta,&\text{若 }a<u(w)-\delta,\\
 a,&\text{若 }a\geq u(w)-\delta,
 \end{cases}
 \tag{14.3}
\end{equation}
其中，$\delta$ 是预先指定的某个身体活动量，例如每周 2 小时。在这一策略下，能够实行的个体须增加 $\delta$ 个单位的身体活动。允许 $\delta$ 成为 $w$ 的函数，还可得到一些值得研究的修正策略，使研究者能够根据健康状况、年龄等协变量指定不同的身体活动增量。

Díaz 和 van der Laan（2012）将这一随机处理策略解释为：改变用于分配暴露水平的概率机制。Haneuse 和 Rotnitzky（2013）指出，这种解释可能并不理想，其论证如下。考虑身体活动的一个新分布，在该分布下，处理按照干预前分布的位置平移版本进行分配。若将随机策略理解为把处理分配机制改为一个位置平移后的分布，那么，原来身体活动水平为 30 分钟的个体，在新策略下可能只获得 10 分钟的处理水平。这可能带来问题，因为这种干预可能降低个体层面的身体活动水平。不过，我们注意到，这两种干预下的暴露总体分布相同，因此会得到完全相同的反事实分布。由此，所采用的解释并不影响因果效应的定义、识别与估计。
\end{bookresult}

\begin{bookresult}{例}{14.2（续）}
一种值得研究的策略为 $d(a,w)=aI(a\leq\delta)+\delta I(a>\delta)$，其中 $I(x)$ 为指示函数：当 $x$ 为真时等于 1，否则等于 0。在这一策略下，要求所有地区的污染水平都不超过 $\delta$。
\end{bookresult}

我们用一个假想世界中结局变量的分布来定义因果量：在这个世界里，分配给个体的是随机策略，而非处理的自然取值。在 NPSEM (14.1) 中，相应的反事实结局定义为 $Y_{d(A,W)}=f_Y(d(A,W),W,U_Y)$。

\subsection{识别}

因果推断路线图的下一步，是识别因果参数。之所以需要识别，是因为反事实变量 $Y_{d(A,W)}$ 通常不可观测。因此，只有当期望 $E(Y_{d(A,W)})$ 能表示为观测数据分布 $P_0$ 的函数时，才可能对它进行估计。具体如下。利用迭代期望法则，我们可以写出
\[
 E(Y_{d(A,W)})=\int_{a\in\mathcal A}\int_{w\in\mathcal W}
 E(Y_{d(a,w)}\mid A=a,W=w)g_0(a\mid w)q_{W,0}(w)\,d\nu(a,w),
\]
其中，$\mathcal A$ 和 $\mathcal W$ 分别为 $A$ 和 $W$ 的分布的支撑集。NPSEM (14.1) 和假设 (14.2) 蕴含：
\begin{enumerate}
 \item $Y_{d(a,w)}\indep A\mid W$；
 \item 在事件 $A=d(a,w)$ 上，$Y_{d(a,w)}=Y$。
\end{enumerate}
于是，期望 $E(Y_{d(A,W)})$ 可由下式识别：
\begin{equation}
\begin{aligned}
 E(Y_{d(A,W)})
 &=\int_{a\in\mathcal A}\int_{w\in\mathcal W}
 E(Y_{d(a,w)}\mid A=d(a,w),W=w)g_0(a\mid w)q_{W,0}(w)\,d\nu(a,w)\\
 &=\int_{a\in\mathcal A}\int_{w\in\mathcal W}
 E(Y\mid A=d(a,w),W=w)g_0(a\mid w)q_{W,0}(w)\,d\nu(a,w)\\
 &=E_{P_0}\{\bar Q(d(A,W),W)\}.
\end{aligned}
\tag{14.4}
\end{equation}
将感兴趣的参数定义为映射 $\Psi:\mathcal M\to\mathbb R$，它把统计模型 $\mathcal M$ 中的一个元素 $P$ 映射为实数 $\Psi(P)$。参数真值是该映射在真实分布 $P_0\in\mathcal M$ 处的取值，记为 $\psi_0=\Psi(P_0)$。因此，感兴趣的统计参数为
\begin{equation}
 \Psi(P)=E_P\{\bar Q(d(A,W),W)\},
 \tag{14.5}
\end{equation}
其中，$\bar Q$ 表示与分布 $P$ 对应的 $Y$ 的条件期望。注意，该参数仅依赖于 $Q=(\bar Q,g,q_W)$。因此，略微放宽记号的使用，我们将交替使用 $\Psi(Q)$ 和 $\Psi(P)$。

\subsection{正值性假设}

在上述识别结果的推导中，我们隐含地假定：对所有 $w\in\mathcal W$，若 $a\in\mathcal A(w)$，则有 $d(a,w)\in\mathcal A(w)$，其中 $\mathcal A(w)$ 表示给定 $W=w$ 时 $A$ 的条件分布支撑集。这一假设通常称为正值性假设（positivity assumption），它保证所考虑的策略得到观测数据的支持。若没有这一假设，式 (14.4) 中的积分可能无法良好定义，因为 $E(Y_{d(a,w)}\mid A=d(a,w),W=w)$ 中的条件集合可能为空。可以认为，这一假设远比静态策略所要求的假设更容易满足；后者要求，在 $W$ 的支撑集上，策略所考虑的所有处理水平都具有正概率。特别是，我们总可以将随机策略定义为满足正值性，式 (14.3) 正是这样做的。

不满足正值性的策略也会使所得因果效应难以解释，因为这样的策略不会在人群中自然出现。为作说明，考虑例 27.1。假定 $w$ 表示冠心病（coronary heart disease，CHD）患者的协变量特征，并且某些 CHD 患者处理的自然取值，恰好等于按其病情推荐的最大身体活动水平 $a$。再假定，其他所有已确诊 CHD 的个体，其 LTPA 值都低于推荐的最大值。此时，策略 $d(a,w)=a+2$ 不满足正值性，因而其效应无法估计。这样的策略也没有多少实际意义，因为强制 CHD 患者实行它并不现实。
